% GENERATED FILE. Edit canonical lessons, then run npm run generate:books.
% canonical-source-hash: fnv1a64:fb3ab2b08de7473b
% canonical-lessons: LA-C467-sutrina, LA-C467-tonstrina, LA-C467-subula, LA-C467-coriarius, LA-C467-pellio

\chapter{Cobbler's shop, Barber's shop, Awl, Tanner, Furrier}
\label{ch:la-cobbler-s-shop-barber-s-shop-awl-tanner-furrier}
\hlchaptermodality{467}

\begin{chapteropening}
\textbf{By the end of this chapter:} \emph{I can talk about small shops and leather work: the cobbler's shop, the barber's, an awl, the tanner and the furrier.}

Complete the last lesson of chapter 467: I can talk about small shops and leather work: the cobbler's shop, the barber's, an awl, the tanner and the furrier.
\end{chapteropening}

\section[sūtrīna, sūtrīnae]{sūtrīna, sūtrīnae — a cobbler's shop}

\label{lesson:LA-C467-sutrina}



\begin{quote}
Before the new one: say the Latin for to value, then the Latin for to be sold.
\end{quote}

\subsection*{You'll want to know: sūtrīna, sūtrīnae}
\textbf{sūtrīna, sūtrīnae} — \textquotedblleft{}a cobbler's shop\textquotedblright{}.

\begin{quote}
Sūtor in sūtrīnā calceōs novōs facit.
\end{quote}

\begin{quote}
\emph{The cobbler makes new shoes in his shop.}
\end{quote}

First declension, feminine: \emph{sūtrīna, sūtrīnae f.}, from \emph{sūtor} 'shoemaker'.

\begin{cousinweb}
\emph{sūtrīna} is built on \emph{sūtor} 'shoemaker' (Lewis \& Short), which comes from \emph{suō} 'to sew'. English \emph{suture} comes from the same verb.
\end{cousinweb}

\begin{grammarlens}[title={What you've built}]
One more word for telling, explaining and describing at B1.
\end{grammarlens}

\subsection*{Your turn}
\begin{itemize}
\raggedright
  \item \emph{Say it:} \emph{sūtrīna, sūtrīnae}
  \item \emph{Say it:} \emph{sūtrīna, sūtrīnae}, once more
  \item \emph{Say it:} a sentence of your own with \emph{sūtrīna, sūtrīnae}
  \item \emph{Recall:} say the Latin for to value, then the Latin for to be sold, then say \emph{sūtrīna, sūtrīnae} again
\end{itemize}

\begin{tcolorbox}[breakable,colback=teal!4,colframe=teal!35!black,title={Before you move on}]
What does sūtrīna, sūtrīnae mean? (\textquotedblleft{}A cobbler's shop\textquotedblright{}.) Say it once more.
\end{tcolorbox}
\section[tōnstrīna, tōnstrīnae]{tōnstrīna, tōnstrīnae — a barber's shop}

\label{lesson:LA-C467-tonstrina}



\begin{quote}
Before the new one: say the Latin for to be sold, then the Latin for a cobbler's shop.
\end{quote}

\subsection*{You'll want to know: tōnstrīna, tōnstrīnae}
\textbf{tōnstrīna, tōnstrīnae} — \textquotedblleft{}a barber's shop\textquotedblright{}.

\begin{quote}
In tōnstrīnā hominēs sedent et fābulās narrant.
\end{quote}

\begin{quote}
\emph{In the barber's shop men sit and tell stories.}
\end{quote}

First declension, feminine: \emph{tōnstrīna, tōnstrīnae f.}, from \emph{tōnsor} 'barber'. A favourite place for gossip.

\begin{cousinweb}
\emph{tōnstrīna} is built on \emph{tōnsor} 'barber' (Lewis \& Short), which comes from \emph{tondeō} 'to shear, to clip'.
\end{cousinweb}

\begin{grammarlens}[title={What you've built}]
One more word for telling, explaining and describing at B1.
\end{grammarlens}

\subsection*{Your turn}
\begin{itemize}
\raggedright
  \item \emph{Say it:} \emph{tōnstrīna, tōnstrīnae}
  \item \emph{Say it:} \emph{tōnstrīna, tōnstrīnae}, once more
  \item \emph{Say it:} a sentence of your own with \emph{tōnstrīna, tōnstrīnae}
  \item \emph{Recall:} say the Latin for to be sold, then the Latin for a cobbler's shop, then say \emph{tōnstrīna, tōnstrīnae} again
\end{itemize}

\begin{tcolorbox}[breakable,colback=teal!4,colframe=teal!35!black,title={Before you move on}]
What does tōnstrīna, tōnstrīnae mean? (\textquotedblleft{}A barber's shop\textquotedblright{}.) Say it once more.
\end{tcolorbox}
\section[sūbula, sūbulae]{sūbula, sūbulae — an awl}

\label{lesson:LA-C467-subula}



\begin{quote}
Before the new one: say the Latin for a cobbler's shop, then the Latin for a barber's shop.
\end{quote}

\subsection*{You'll want to know: sūbula, sūbulae}
\textbf{sūbula, sūbulae} — \textquotedblleft{}an awl\textquotedblright{}.

\begin{quote}
Sūtor sūbulā et fīlō calceōs facit.
\end{quote}

\begin{quote}
\emph{The cobbler makes shoes with an awl and thread.}
\end{quote}

First declension, feminine: \emph{sūbula, sūbulae f.} A pointed tool for piercing leather before sewing.

\begin{cousinweb}
\emph{sūbula} is inherited from PIE \emph{syuh\textsubscript{1}-d\textsuperscript{h}leh\textsubscript{2}} (EtymDB), built on the root of \emph{suō} 'to sew': a tool for sewing leather.
\end{cousinweb}

\begin{grammarlens}[title={What you've built}]
One more word for telling, explaining and describing at B1.
\end{grammarlens}

\subsection*{Your turn}
\begin{itemize}
\raggedright
  \item \emph{Say it:} \emph{sūbula, sūbulae}
  \item \emph{Say it:} \emph{sūbula, sūbulae}, once more
  \item \emph{Say it:} a sentence of your own with \emph{sūbula, sūbulae}
  \item \emph{Recall:} say the Latin for a cobbler's shop, then the Latin for a barber's shop, then say \emph{sūbula, sūbulae} again
\end{itemize}

\begin{tcolorbox}[breakable,colback=teal!4,colframe=teal!35!black,title={Before you move on}]
What does sūbula, sūbulae mean? (\textquotedblleft{}An awl\textquotedblright{}.) Say it once more.
\end{tcolorbox}
\section[coriārius, coriāriī]{coriārius, coriāriī — a tanner, a currier}

\label{lesson:LA-C467-coriarius}



\begin{quote}
Before the new one: say the Latin for a barber's shop, then the Latin for an awl.
\end{quote}

\subsection*{You'll want to know: coriārius, coriāriī}
\textbf{coriārius, coriāriī} — \textquotedblleft{}a tanner, a currier\textquotedblright{}.

\begin{quote}
Coriārius coria bovum in aquā lavat.
\end{quote}

\begin{quote}
\emph{The tanner washes the hides of the oxen in water.}
\end{quote}

Second declension, masculine: \emph{coriārius, coriāriī m.}, from \emph{corium} 'hide, leather'. Tanneries stank and stood at the edge of town. English \emph{currier} comes from it.

\begin{cousinweb}
\emph{coriārius} is built on \emph{corium} 'hide, leather' with the suffix \emph{-ārius} (Lewis \& Short). English \emph{currier}, a leather dresser, comes from it through French.
\end{cousinweb}

\begin{grammarlens}[title={What you've built}]
One more word for telling, explaining and describing at B1.
\end{grammarlens}

\subsection*{Your turn}
\begin{itemize}
\raggedright
  \item \emph{Say it:} \emph{coriārius, coriāriī}
  \item \emph{Say it:} \emph{coriārius, coriāriī}, once more
  \item \emph{Say it:} a sentence of your own with \emph{coriārius, coriāriī}
  \item \emph{Recall:} say the Latin for a barber's shop, then the Latin for an awl, then say \emph{coriārius, coriāriī} again
\end{itemize}

\begin{tcolorbox}[breakable,colback=teal!4,colframe=teal!35!black,title={Before you move on}]
What does coriārius, coriāriī mean? (\textquotedblleft{}A tanner, a currier\textquotedblright{}.) Say it once more.
\end{tcolorbox}
\section[pelliō, pelliōnis]{pelliō, pelliōnis — a furrier, a worker in skins}

\label{lesson:LA-C467-pellio}



\begin{quote}
Before the new one: say the Latin for an awl, then the Latin for a tanner.
\end{quote}

\subsection*{You'll want to know: pelliō, pelliōnis}
\textbf{pelliō, pelliōnis} — \textquotedblleft{}a furrier, a worker in skins\textquotedblright{}.

\begin{quote}
Pelliō vestem calidam facit.
\end{quote}

\begin{quote}
\emph{The furrier makes a warm garment.}
\end{quote}

Third declension, masculine: \emph{pelliō, pelliōnis m.} He made and sold warm clothes of skins.

\begin{cousinweb}
\emph{pelliō} is built on \emph{pellis} 'skin, hide' with the suffix \emph{-iō} for a person who works with something (Lewis \& Short). English \emph{pelt} is often linked with \emph{pellis}, but the link is not settled.
\end{cousinweb}

\begin{grammarlens}[title={What you've built}]
One more word for telling, explaining and describing at B1.
\end{grammarlens}

\subsection*{Your turn}
\begin{itemize}
\raggedright
  \item \emph{Say it:} \emph{pelliō, pelliōnis}
  \item \emph{Say it:} \emph{pelliō, pelliōnis}, once more
  \item \emph{Say it:} a sentence of your own with \emph{pelliō, pelliōnis}
  \item \emph{Recall:} say the Latin for an awl, then the Latin for a tanner, then say \emph{pelliō, pelliōnis} again
\end{itemize}

\begin{tcolorbox}[breakable,colback=teal!4,colframe=teal!35!black,title={Before you move on}]
What does pelliō, pelliōnis mean? (\textquotedblleft{}A furrier, a worker in skins\textquotedblright{}.) Say it once more.
\end{tcolorbox}
